\documentclass[11pt]{article}
\usepackage[letterpaper,margin=0.72in,top=0.82in,bottom=0.72in,headheight=22pt]{geometry}
\usepackage[T1]{fontenc}
\usepackage{lmodern}
\usepackage{microtype}
\usepackage[table]{xcolor}
\usepackage{amsmath,amssymb,mathtools}
\usepackage{array,longtable,booktabs,tabularx}
\usepackage[inline]{enumitem}
\usepackage[most]{tcolorbox}
\usepackage{fancyhdr}
\usepackage{hyperref}

\definecolor{Primary}{HTML}{17324D}
\definecolor{Accent}{HTML}{2C7A7B}
\definecolor{CueTint}{HTML}{EAF1F5}
\definecolor{NoteTint}{HTML}{F8FAFB}
\definecolor{Rule}{HTML}{B8C6CF}
\hypersetup{hidelinks,pdftitle={Chapter 5: Introduction to Suffix Trees},pdfsubject={Cornell notes}}
\setlist[itemize]{leftmargin=1.25em,itemsep=1pt,topsep=1.5pt,parsep=0pt}
\setlength{\parindent}{0pt}
\setlength{\parskip}{3pt}
\renewcommand{\arraystretch}{1.16}
\emergencystretch=2em

\pagestyle{fancy}
\fancyhf{}
\fancyhead[L]{\sffamily\small\color{Primary} Chapter 5}
\fancyhead[R]{\sffamily\small\color{Primary} Suffix-tree basics}
\fancyfoot[C]{\sffamily\small\color{Primary}\thepage}
\renewcommand{\headrulewidth}{0.5pt}
\renewcommand{\headrule}{\hbox to\headwidth{\color{Accent}\leaders\hrule height \headrulewidth\hfill}}

\newcolumntype{C}{>{\raggedright\arraybackslash}p{0.275\textwidth}}
\newcolumntype{N}{>{\raggedright\arraybackslash}p{0.655\textwidth}}
\newcommand{\cnrow}[2]{%
  \cellcolor{CueTint}\textbf{\sffamily\color{Primary}#1} &
  \cellcolor{NoteTint}#2 \\[4pt]\arrayrulecolor{Rule}\hline}

\begin{document}

\begin{tcolorbox}[
  enhanced,breakable,colback=Primary!4,colframe=Primary,boxrule=0.9pt,
  arc=2mm,left=3mm,right=3mm,top=2.5mm,bottom=2.5mm]
{\sffamily\color{Accent}\bfseries CHAPTER 5}\\[-1pt]
{\sffamily\LARGE\bfseries\color{Primary} Introduction to Suffix Trees}

\vspace{2pt}
\textbf{Purpose.} Define the compressed index of all text suffixes and show why it supports fast substring queries while occupying linear space.

\textbf{Learning objectives}
\begin{itemize}
  \item Define explicit and implicit nodes, edge labels, path labels, and suffix leaves.
  \item Explain why a unique terminator makes every suffix end at a leaf.
  \item Prove the linear-size bound for a compact suffix tree.
  \item Contrast a naive construction with later linear-time algorithms.
\end{itemize}

\textbf{Prerequisites.} Tries, rooted trees, substrings and suffixes, and compact edge-label representations.
\end{tcolorbox}

\begin{longtable}{@{}C!{\color{Accent}\vrule width 0.8pt}N@{}}
\rowcolor{Primary}
\color{white}\sffamily\bfseries Cue / question &
\color{white}\sffamily\bfseries Notes \\ \hline
\endfirsthead
\rowcolor{Primary}
\color{white}\sffamily\bfseries Cue / question &
\color{white}\sffamily\bfseries Notes (continued) \\ \hline
\endhead
\cnrow{\S5.1: Why is the structure important?}{A suffix tree preprocesses a text once so that a later pattern can be located by following one path. It combines the query power of a trie containing every suffix with path compression that removes unary internal nodes.}
\cnrow{\S5.2: What is a suffix tree?}{For $S\$$, where $\$$ is unique and lexicographically distinct if ordering is needed, build the compact trie of all suffixes. Each leaf identifies a suffix start. The concatenation of labels from the root to a node is its path label; edges leaving one node begin with distinct characters.}
\cnrow{What are explicit and implicit locations?}{Explicit nodes are stored branching points, the root, and leaves. A point inside a compressed edge is an implicit node or locus. A pattern may end at either kind of locus; its descendant leaves are exactly the occurrence positions.}
\cnrow{Why is the tree linear in size?}{There are exactly $n+1$ leaves for the suffixes of $S\$$. Every nonroot internal explicit node has at least two children, so there are fewer internal nodes than leaves. Store each edge label as a pair of text indices rather than copying characters; the full representation is $O(n)$ space.}
\cnrow{\S5.3: How does a query work?}{Starting at the root, choose an outgoing edge by the next pattern character and compare along its text interval. Consuming the entire pattern identifies a locus; enumerate descendant leaf labels to report occurrences. Search costs $O(|P|+\mathrm{occ})$ with constant-time or suitably bounded edge selection.}
\cnrow{What does a small example reveal?}{For $S=\texttt{banana}\$$, repeated prefixes of suffixes such as \texttt{ana} share a path. Leaves below the \texttt{ana} locus identify starts $2$ and $4$ under one-based indexing. Compression preserves branching decisions while removing repeated storage.}
\cnrow{\S5.4: Why is naive construction expensive?}{Insert all suffixes into a trie and compress unary paths. The suffixes have total length $\Theta(n^2)$, so direct insertion can use quadratic time and space before compression. A linear-size final structure does not by itself imply a linear-time construction.}
\cnrow{What implementation details are already visible?}{Use a unique terminator, store edge intervals into the original text, attach suffix-start labels to leaves, and define the child dictionary according to alphabet size. Queries must detect a mismatch inside an edge, not only at explicit nodes.}
\end{longtable}


\begin{tcolorbox}[enhanced,breakable,title={Chapter synthesis},
  colback=Accent!5,colframe=Accent,fonttitle=\small\sffamily\bfseries,fontupper=\footnotesize,
  boxsep=0.6mm,left=1.5mm,right=1.5mm,top=1mm,bottom=1mm,before skip=3pt,after skip=3pt]
A suffix tree is a compact trie of all terminated suffixes. Shared prefixes become common paths, unary paths become text-indexed edges, and descendant leaves encode occurrence positions. The result is a linear-space text index with pattern-time queries, although obtaining it naively is quadratic.
\end{tcolorbox}

\begin{tcolorbox}[enhanced,breakable,title={Key takeaways},
  colback=white,colframe=Primary,fonttitle=\small\sffamily\bfseries,fontupper=\footnotesize,
  boxsep=0.6mm,left=1.5mm,right=1.5mm,top=1mm,bottom=1mm,before skip=3pt,after skip=3pt]
\begin{itemize}
  \item The terminator makes every suffix explicit and unique.
  \item Edge labels should reference the text rather than copy substrings.
  \item A pattern locus may lie inside an edge.
  \item Occurrence positions are stored at descendant leaves.
  \item Branching degree yields the linear node bound.
  \item Output enumeration must appear in the query complexity.
\end{itemize}
\end{tcolorbox}

\begin{tcolorbox}[enhanced,breakable,title={Algorithm selection / when to use this},
  colback=Primary!3,colframe=Primary,fonttitle=\small\sffamily\bfseries,fontupper=\footnotesize,
  boxsep=0.6mm,left=1.5mm,right=1.5mm,top=1mm,bottom=1mm,before skip=3pt,after skip=3pt]
\begin{itemize}
  \item Use a suffix tree when one text supports many substring or structural queries.
  \item Use a simpler pattern-preprocessing matcher for one-off searches when index construction is unnecessary.
  \item Use a suffix array when space constants and cache behavior dominate and its query tradeoffs are acceptable.
\end{itemize}
\end{tcolorbox}

\begin{tcolorbox}[enhanced,breakable,title={Self-test questions},
  colback=white,colframe=Accent,fonttitle=\small\sffamily\bfseries,fontupper=\footnotesize,
  boxsep=0.6mm,left=1.5mm,right=1.5mm,top=1mm,bottom=1mm,before skip=3pt,after skip=3pt]
\begin{itemize}
  \item Why does omitting the terminator permit a suffix to end inside an edge?
  \item Prove that the number of branching internal nodes is less than the number of leaves.
  \item How is a pattern ending mid-edge represented?
  \item Why does naive suffix insertion take quadratic work?
  \item What must be included in the time to report all occurrences?
\end{itemize}
\end{tcolorbox}

{\footnotesize
\textbf{Terms and notation to remember.} suffix tree, compact trie, terminator, path label, explicit node, implicit node, locus, leaf label, edge interval.

\textbf{Connections.} The next chapter constructs the same linear-size structure in linear time; later chapters augment it with colors, depths, suffix arrays, and lowest-common-ancestor queries.
}

\end{document}
